\documentclass[11pt]{article}

\usepackage{amsmath,amssymb,graphicx}
\usepackage{hyperref}
\usepackage{geometry}
\geometry{margin=1in}

\title{On the Effective Single-Harmonic Structure of Local Energy Landscapes in Molecular VQE}

\author{Your Name}

\date{}

\begin{document}

\maketitle

% ============================================================
\begin{abstract}

Variational Quantum Algorithms (VQAs) are widely used for estimating ground-state energies of molecular systems, yet their optimization remains challenging due to the non-convex structure of the energy landscape. It is known that the energy of parameterized quantum circuits admits a finite Fourier representation and that this structure can be exploited for coordinate-wise optimization.

In this work, we focus on a narrower empirical question: what is the effective spectral structure of VQE energy landscapes in small molecular problems? We show that, for H$_2$, LiH, and BeH$_2$, the energy along one-dimensional directions in parameter space is often strongly dominated by the first Fourier harmonic in the vicinity of the optimum, leading to near-zero spectral entropy.

Motivated by this observation, we propose a simple first-harmonic surrogate that approximates the local landscape and provides an analytical estimate of the minimum along a given direction. We demonstrate that this approach achieves a favorable cost-accuracy trade-off compared to higher-order trigonometric reconstruction, requiring fewer evaluations while maintaining good optimization performance.

Our results suggest that, in practical molecular VQE settings, exact reconstruction of the full trigonometric structure is often unnecessary, and that low-cost approximations based on dominant spectral components can be effective.

\end{abstract}

% ============================================================
\section{Introduction}

Variational Quantum Algorithms (VQAs) have emerged as a leading approach for estimating ground-state energies of molecular systems on near-term quantum devices. In particular, the Variational Quantum Eigensolver (VQE) combines quantum circuit evaluations with classical optimization.

Despite its success, VQE optimization remains challenging due to the non-convex and potentially complex structure of the energy landscape. A growing body of work has studied this landscape from both geometric and functional perspectives, including convergence guarantees and Fourier representations of parameterized quantum circuits.

It is known that, for circuits composed of parametrized rotations, the expectation value of observables can be expressed as a finite Fourier series in each parameter. This structure has been exploited to derive coordinate-wise optimization methods that reconstruct and minimize these trigonometric functions exactly.

In this work, we adopt a different perspective. Instead of focusing on exact reconstruction, we investigate the \emph{effective} spectral structure of VQE landscapes in practical molecular settings. Specifically, we ask whether the full trigonometric structure is necessary in practice, or whether simpler approximations suffice.

Our main observation is that, for small molecular systems and near good solutions, the energy along one-dimensional directions is often dominated by a single Fourier component. This allows the use of a much simpler surrogate model that approximates the minimum with significantly reduced cost.

The contributions of this work are:
\begin{itemize}
    \item An empirical characterization of the local Fourier structure of VQE landscapes in molecular systems.
    \item Evidence that first-harmonic dominance occurs consistently near optimal solutions.
    \item A simple first-harmonic surrogate for estimating coordinate-wise minima.
    \item An analysis of the cost-accuracy trade-off compared to higher-order reconstruction.
\end{itemize}

% ============================================================
\section{Background}

\subsection{Variational Quantum Eigensolver}

In VQE, the energy is given by
\begin{equation}
E(\boldsymbol{\theta}) = \langle \psi(\boldsymbol{\theta}) | H | \psi(\boldsymbol{\theta}) \rangle,
\end{equation}
where $|\psi(\boldsymbol{\theta})\rangle = U(\boldsymbol{\theta}) |\psi_0\rangle$ is a parameterized quantum state.

\subsection{Fourier Structure of Parameterized Circuits}

For parameterized circuits composed of rotations generated by Hermitian operators with discrete spectra, the energy as a function of a single parameter $\theta$ can be written as a finite Fourier series:
\begin{equation}
E(\theta) = \sum_{k} a_k \cos(k\theta) + b_k \sin(k\theta) + c.
\end{equation}

Previous works have exploited this structure to reconstruct the function exactly and compute coordinate-wise optima analytically.

% ============================================================
\section{Local Harmonic Structure}

While the global structure of $E(\theta)$ may involve multiple harmonics, we focus on its behavior near a local optimum.

Empirically, we observe that near optimal parameters $\boldsymbol{\theta}^*$, the energy along directions of the form
\begin{equation}
\boldsymbol{\theta}(\phi) = \boldsymbol{\theta}^* + \phi \mathbf{v}
\end{equation}
is well approximated by a single harmonic:
\begin{equation}
E(\phi) \approx A \cos(\phi - \delta) + C.
\end{equation}

This implies that higher-order Fourier components are effectively suppressed near the optimum, leading to a low-entropy spectrum.

% ============================================================
\section{Methodology}

We evaluate this phenomenon for the molecules H$_2$, LiH, and BeH$_2$, using parity mapping with $\mathbb{Z}_2$ tapering.

For each system, we:
\begin{enumerate}
    \item Run VQE to obtain an approximate optimum $\boldsymbol{\theta}^*$.
    \item Sample directions $\mathbf{v}$ in parameter space.
    \item Evaluate $E(\phi)$ along $\boldsymbol{\theta}^* + \phi \mathbf{v}$.
    \item Fit a first-harmonic model:
    \begin{equation}
    E(\phi) \approx \frac{a_0}{2} + a_1 \cos\phi + b_1 \sin\phi.
    \end{equation}
\end{enumerate}

The estimated minimum is given by
\begin{equation}
\phi^* = \mathrm{atan2}(b_1, a_1) + \pi.
\end{equation}

We compare this approximation with higher-order Fourier reconstruction and direct sampling.

% ============================================================
\section{Results}

\subsection{Spectral Structure}

Across all tested systems, we observe that the first harmonic dominates the spectrum near optimal solutions, with spectral entropy close to zero.

\subsection{Local vs Global Behavior}

Away from the optimum, the spectrum becomes richer and higher-order components become relevant. This confirms that the monofrequency structure is a local phenomenon.

\subsection{Cost-Accuracy Trade-off}

We compare:
\begin{itemize}
    \item First-harmonic approximation (K=1)
    \item Higher-order Fourier reconstruction
\end{itemize}

Results show that the first-harmonic model provides a good approximation of the minimum while requiring significantly fewer evaluations.

% ============================================================
\section{Discussion}

Our results suggest that, in small molecular VQE instances, the local energy landscape is effectively low-dimensional and dominated by a single harmonic component.

This explains why simple optimization strategies often perform well in practice and indicates that exact reconstruction of the full trigonometric structure may be unnecessary in many cases.

The main contribution is therefore not a new optimization principle, but an empirical observation that enables a cheaper approximation with favorable performance.

% ============================================================
\section{Conclusion}

We have shown that VQE energy landscapes for small molecular systems exhibit strong first-harmonic dominance near optimal solutions.

This allows the use of a simple surrogate model that approximates coordinate-wise minima with reduced computational cost.

Future work includes extending this approach to multi-parameter updates and larger systems.

% ============================================================

\bibliographystyle{plain}
\bibliography{refs}

\end{document}